% =============================================================
% 4.5 การออกแบบเชิงการศึกษาและการประเมินผล
% โครงงาน: การออกแบบและพัฒนาเกมเดินเรือด้วยระบบสั่งการด้วยเสียงเพื่อส่งเสริมทักษะการสื่อสาร
% =============================================================

\section{การออกแบบเชิงการศึกษาและการประเมินผล}

เนื่องจากโครงงานนี้เป็นเกมเพื่อการเรียนรู้ (Serious Game) การออกแบบจึงบูรณาการการเรียนรู้โดยใช้เกมเป็นฐาน (Digital Game-Based Learning) เข้ากับทฤษฎีการเรียนรู้ภาษา เพื่อให้ผู้เล่นได้พัฒนาทักษะเชิงปฏิบัติ~\cite{sadigzade2025} โดยแบ่งออกเป็น 2 ส่วน ดังนี้

\subsection{การออกแบบเชิงการศึกษา (Educational Design)}

การออกแบบกลไกเพื่อส่งเสริมการเรียนรู้ประกอบด้วยแนวคิดหลัก ดังนี้

\begin{enumerate}[label=\arabic*)]
    \item \textbf{การเรียนรู้ตามสถานการณ์ (Situated Learning):} ผู้เล่นไม่ได้เรียนรู้คำศัพท์จากการท่องจำ แต่เรียนรู้ผ่านการจำลองการสื่อสารทางวิทยุในการเดินเรือ เช่น การสั่งหางเสือหรือการปรับความเร็ว เพื่อให้เข้าใจภาษาในบริบทที่ใช้จริง~\cite{lave1991} ซึ่งสอดคล้องกับแนวทางการเรียนรู้โดยใช้เกมเป็นฐาน~\cite{sadigzade2025}
    \item \textbf{การลดตัวกรองทางอารมณ์ (Affective Filter Hypothesis):} สมมติฐานนี้เสนอว่าความวิตกกังวลขัดขวางการรับภาษา~\cite{krashen1982} การให้ผู้เล่นฝึกพูดสั่งการกับลูกเรือปัญญาประดิษฐ์แทนการสนทนากับมนุษย์ จึงออกแบบมาเพื่อสร้างพื้นที่ปลอดภัย (Safe Space) ลดความกลัวความผิดพลาด และบรรเทาความวิตกกังวลในการใช้ภาษาต่างประเทศ (Language Anxiety)
    \item \textbf{การเรียนรู้จากข้อผิดพลาด (Learning from Failure):} เมื่อผู้เล่นออกเสียงไม่ชัดเจน ระบบไม่ลงโทษด้วยการตัดจบทันที แต่ลูกเรือจะขอให้พูดซ้ำ และเมื่อผู้เล่นใช้รูปประโยคที่ไม่ใช่มาตรฐาน ลูกเรือจะทวนกลับด้วยรูปประโยคมาตรฐาน (Recast) ทั้งสองแบบเป็นข้อมูลป้อนกลับ (Feedback) ที่กระตุ้นให้ผู้เล่นแก้ไขการพูดด้วยตนเอง~\cite{zargham2024}
\end{enumerate}

\subsection{แนวทางการประเมินผล (Evaluation Approach)}

ผู้จัดทำออกแบบการประเมินผลให้ครอบคลุมทั้งประสิทธิภาพของเทคโนโลยีและผลด้านการเรียนรู้ของผู้เล่น ดังนี้

\begin{itemize}
    \item \textbf{การประเมินผลภายในเกม (In-game Evaluation):} ประเมินจากร้อยละของการสื่อสารที่ลูกเรือเข้าใจ จำนวนครั้งที่ลูกเรือขอให้พูดซ้ำ ความรวดเร็วในการสั่งการหลังได้รับคำเตือน จำนวนครั้งที่ชน และจำนวนคำสั่งฉุกเฉินที่ทันเวลา
    \item \textbf{การประเมินประสิทธิภาพทางเทคนิค (Technical Evaluation):} ประเมินความแม่นยำของการแปลงเสียงเป็นข้อความโดยให้ผู้เข้าร่วมแต่ละคนอ่านรายการคำสั่งมาตรฐาน 10 คำสั่ง (คำสั่งเป้าหมาย) ระหว่างการเล่น แล้วคำนวณร้อยละของคำสั่งเป้าหมายที่ระบบถอดความและแยกคำสั่งได้ถูกต้อง โดยเทียบรายการคำสั่งเป้าหมายกับคอลัมน์ข้อความที่ได้ยินและคำสั่งที่ระบบแยกได้ในแฟ้ม turns.csv ส่วนการเล่นอิสระรายงานเฉพาะร้อยละของการสื่อสารที่ลูกเรือเข้าใจ นอกจากนี้วัดความหน่วง (Latency) ตั้งแต่ผู้เล่นปล่อยปุ่มวิทยุจนเรือเริ่มตอบสนอง
    \item \textbf{การประเมินความวิตกกังวลในการพูดภาษาอังกฤษ (Speaking Anxiety Evaluation):} ผู้เล่นตอบแบบวัดความวิตกกังวลในการพูดภาษาอังกฤษ 12 ข้อ ที่ดัดแปลงจาก FLCAS~\cite{horwitz1986} ทั้งก่อนและหลังเล่นเกม เพื่อเปรียบเทียบว่าความวิตกกังวลลดลงหรือไม่
    \item \textbf{การประเมินประสบการณ์ผู้เล่น (User Experience Evaluation):} ประเมินผ่านแบบสอบถามหลังการเล่น 10 ข้อ ครอบคลุมความเข้าใจง่าย ความสมจริง การตอบกลับของลูกเรือ ประโยชน์ของคำแนะนำประโยคและการทวนด้วยรูปประโยคมาตรฐาน ความรู้สึกปลอดภัยในการพูด และความสนุก โดยแบบวัดและแบบสอบถามทั้งหมดแสดงในภาคผนวก ก และวิธีวิเคราะห์ผลอธิบายในบทที่ 3
\end{itemize}